\begin{chapintro}
検査の研究は感度と特異度という 2 つの結果を持ち、カットオフや対象の違いによって 2 つが連動して動きます。そのため、感度と特異度を別々に平均する方法は勧められず、両者を同時に扱う二変量モデルか HSROC モデルを使います。要約点を出すか SROC 曲線を出すかはカットオフがそろっているかで決まり、異質性と小規模研究効果の調べ方も、介入のメタ解析とは違います。二変量モデルの仕組み、HSROC モデル、異質性と小規模研究効果の細部は発展の節に置いてあり、初読では飛ばせます。
\end{chapintro}

\section{2 つの結果が連動して動く}
\label{sec:10-1}

第 1 部の介入のメタ解析では、試験ごとにリスク比を 1 つ計算し、その重み付き平均を求めました。検査の研究では、研究ごとに感度と特異度の 2 つの数があります。2 つを別々に扱ってよいなら、第 1 部の方法を 2 回使えば済みますが、そうはいきません。

\kkref{sec:09-5}で見たように、同じ検査でもカットオフを下げれば感度が上がって特異度が下がり、上げれば逆になります。メタ解析に集めた研究が別々のカットオフを使っていると、研究ごとの感度と特異度は 1 本の ROC 曲線に沿って散らばります。カットオフを明示していなくても、測定試薬の違いや、判定する人の基準の厳しさの違いが、実質的にカットオフをずらすことがあります。この現象を閾値効果と呼びます。

例題の架空の 10 研究（カットオフは 0.25 ng/mL が 2 研究、0.5 ng/mL が 6 研究、1 ng/mL が 2 研究）を、横軸に偽陽性割合、縦軸に感度をとった ROC 平面に置いたのが図 10.1 です。点の大きさは研究の人数を、色の濃さはカットオフを表します。カットオフ 0.25 ng/mL の 2 研究（Nguyen 2014、Tanaka 2023）は右上に寄り、感度は 0.80 以上ですが特異度は 0.71 と 0.54 にとどまります。カットオフ 1 ng/mL の Sato 2021 は特異度 0.91 で左側にあります。研究どうしの感度と特異度の違いの一部は、カットオフの違いで説明できそうです。

\kkfig[.62\linewidth]{p2_roc_points.pdf}{図 10.1　架空の PCT 研究 10 本を ROC 平面に置いた図。色が濃いほどカットオフが高い}

同じ図には、カットオフでは説明できない散らばりも見えます。カットオフ 0.5 ng/mL の 6 研究（Kato 2010、Lopez 2011、Mori 2013、Park 2016、Qureshi 2018、Rossi 2019）の感度は 0.56 から 0.88 までばらついています。対象患者の重症度、採血のタイミング、測定法の違い、それに偶然が重なった結果です。検査のメタ解析では、研究間のばらつきが大きいほうが普通で、それを前提にしたモデルが必要になります。

\section{感度と特異度を別々に統合すると何が起きるか}
\label{sec:10-2}

最も素朴な方法は、感度だけを集めてランダム効果モデルで統合し、特異度だけを集めて別に統合することです。ランダム効果モデルは\kkref{sec:05-4}で説明しました。カットオフ 0.5 ng/mL の 6 研究では、感度 0.746、特異度 0.757 になります。実データの AUDIT-C（14 研究）でも、別々に統合した感度 0.930、特異度 0.783 と、二変量モデルの感度 0.925、特異度 0.782 は近い値になります。

点推定が近くても、別々に統合する方法には限界があります。AUDIT-C の二変量モデルは、感度と特異度の相関を −0.818 と推定しました。別々に統合すると、この相関は分からないままです。1 つ目に、相関を無視しているので、2 つを組み合わせた点の不確かさを正しく表せません。信頼区間を 2 本別々に描くと長方形の範囲になりますが、実際の不確かさは傾いた楕円の形をしています。2 つ目に、予測領域や SROC 曲線のように、感度と特異度の関係にもとづく図が描けません。3 つ目に、閾値効果が強いデータでは、別々に平均した点が研究の散らばる帯から外れた位置に来ることがあり、そのような感度と特異度の組み合わせを示す研究はどこにもない、ということが起きます。点推定が近いことと、二変量モデルが要らないことは同じではありません。

\begin{hattenbox}
古い検査のメタ解析では、診断オッズ比だけを統合する方法や、Moses と Littenberg の方法で SROC 曲線を引く方法も使われてきました。後者は、研究ごとの診断オッズ比の対数を、感度と偽陽性割合から作った閾値の代わりの量に回帰させて曲線を描きます。簡便ですが、研究間のばらつきを組み込まない固定効果の方法です。Cochrane DTA ハンドブックは、次に説明する二変量モデルと HSROC モデルが、これらの方法に取って代わったと述べています（Macaskill 2023、9.4）。古い論文を読むときは、どの方法で統合したかをまず確かめてください。
\end{hattenbox}

\section{感度と特異度を logit の目盛りに直す}
\label{sec:10-3}

二変量モデルは、感度と特異度を割合のまま扱わず、logit という目盛りに直して扱います。割合 p の logit は log(p ÷ (1 − p)) で、オッズ（\kkref{sec:05-1-4}）の対数です。0 から 1 の範囲に閉じた割合を、端のない目盛りに広げます。第 1 部でリスク比の対数をとって統合したのと同じ発想です\kksee{sec:05-1-2}。表 10.1 に、いくつかの割合の logit を示します。

\begin{kkcenter}
\begin{tabular}{rrr}
\toprule
割合 p & オッズ & logit \\
\midrule
0.50 & 1 & 0 \\
0.75 & 3 & 1.10 \\
0.90 & 9 & 2.20 \\
0.95 & 19 & 2.94 \\
0.99 & 99 & 4.60 \\
\bottomrule
\end{tabular}
\kkcap{表 10.1　割合とその logit。R の qlogis() で計算した}
\end{kkcenter}

割合の差が同じ 0.04 でも、0.50 と 0.54 では logit の差は約 0.16、0.95 と 0.99 では約 1.65 になります。logit の目盛りを使う理由は 2 つあります。1 つ目は、割合のまま平均や区間を考えると、端（0 や 1 の近く）で区間が 0〜1 をはみ出すことがあるからです。2 つ目は、端では割合の差が小さく見えても意味が大きいからです。特異度 0.95 と 0.99 は、偽陽性が 100 人に 5 人か 1 人かという 5 倍の違いにあたります。R の plogis() は logit から元の割合に戻し、plogis(2.94) は約 0.95 を返します。

\section{二変量モデルの考え方}
\label{sec:10-4}

二変量モデルは、感度と特異度を一緒に扱うモデルです（Reitsma 2005）。研究ごとに真の感度と真の特異度が違うと考え、その散らばりを logit の目盛りで表します。この散らばりは、第 1 部のランダム効果モデルの研究間分散 τ² に当たります\kksee{sec:05-4}。感度と特異度を別々でなく同時に扱うので、2 つの間の相関（感度の高い研究ほど特異度が低くなる傾向）も推定できます。結果は、要約点（平均の感度と特異度）と、そのまわりの楕円（信頼領域と予測領域）で示します。楕円の傾きが相関に、楕円の大きさが研究間のばらつきに当たります。

相関は、−1 なら研究の点が完全に一直線に並び、0 なら感度と特異度に関係がないことを表します。研究が少ないと相関がうまく推定できず、−1 や 1 に張り付くことがあります。そのときは相関を 0 と置いた簡単なモデルに切り替えます（\kkref{sec:13-4}）。

\subsection{\hatten 二変量モデルの仕組み}
\label{sec:10-4-1}

二変量モデルは、感度と特異度を 2 つの階で同時に扱います。図 10.2 にその構造を示します。

\begin{kkcenter}
\begin{tikzpicture}[node distance=7mm and 6mm]
  \node[kbox, align=flush center, text width=96mm] (top) {研究間の階：研究ごとの「真の」logit 感度と logit 特異度は、\\平均（要約点）のまわりに 2 次元の正規分布で散らばる\\（2 つの平均、2 つの分散、1 つの相関）};
  \node[kplain, align=flush center, text width=25mm, below=of top, xshift=-33mm] (s1) {研究 1 の真の\\感度と特異度};
  \node[kplain, align=flush center, text width=25mm, below=of top] (s2) {研究 2 の真の\\感度と特異度};
  \node[kplain, align=flush center, text width=25mm, below=of top, xshift=33mm] (s3) {研究 k の真の\\感度と特異度};
  \node[kplain, align=flush center, text width=27mm, below=of s1, fill=kkpale] (o1) {観測した TP、FN\\と TN、FP};
  \node[kplain, align=flush center, text width=27mm, below=of s2, fill=kkpale] (o2) {観測した TP、FN\\と TN、FP};
  \node[kplain, align=flush center, text width=27mm, below=of s3, fill=kkpale] (o3) {観測した TP、FN\\と TN、FP};
  \draw[kar] (top) -- (s1); \draw[kar] (top) -- (s2); \draw[kar] (top) -- (s3);
  \draw[kar] (s1) -- (o1); \draw[kar] (s2) -- (o2); \draw[kar] (s3) -- (o3);
  \node[font=\footnotesize, anchor=west, text=kkdark] at ($(s3.east)!.5!(o3.east)+(1mm,0)$) {研究内の階};
  \node[font=\scriptsize, align=flush center, below=2mm of o2, text width=100mm] {研究内の階：真の感度のもとで、病気のある人のうち何人が陽性になるかは二項分布に従う\\（特異度も同じ）};
\end{tikzpicture}
\kkcap{図 10.2　二変量モデルの 2 つの階}
\end{kkcenter}

下の階（研究内の階）は、各研究の中で起きる偶然の誤差を表します。ある研究の真の感度が 0.75 でも、菌血症の患者 45 人のうち陽性になる人数は偶然で上下し、ちょうど 34 人になるとは限りません。人数の少ない研究ほど、観測した感度は真の値から大きく外れやすくなります。

上の階（研究間の階）は、研究ごとに真の感度と特異度そのものが違うことを表します。カットオフ、対象患者、測定法の違いによって、研究 1 の真の感度と研究 2 の真の感度は異なります。二変量モデルでは、感度と特異度を logit の目盛りに直したうえで、その組が 2 次元の正規分布に従って散らばると考えます。

2 次元の正規分布は、5 つの数で決まります。logit 感度の平均と logit 特異度の平均（この 2 つを元の目盛りに戻したものが要約点です）、それぞれの研究間のばらつき（分散）、そして 2 つの相関です。閾値効果があると、感度が高い研究ほど特異度が低くなるので、logit 感度と logit 特異度の相関は負になります。言い換えると、感度と偽陽性割合の相関は正になります。AUDIT-C では、この相関は −0.818 と推定されました。一方、例題のカットオフ 0.5 ng/mL の 6 研究では、\kkref{sec:13-3}で見るように相関が −1 の境界に張り付き、推定できませんでした。研究が少ないと、相関はこのように決まりにくくなります。

研究内の階の扱い方には 2 通りがあります。Reitsma らの元の方法は、研究ごとの logit 感度と logit 偽陽性割合を正規分布で近似し、その分散を人数から計算して与えます。この近似は、人数の少ない研究や、感度や特異度が 1 に近い研究で悪くなり、0 の升目があれば 0.5 を足す補正も要ります。Chu と Cole（2006）が示した方法は、研究内の人数を二項分布のまま扱う一般化線形混合モデルで、0 の升目がある研究でも補正なしに扱えます。Cochrane DTA ハンドブックの第 10 章はこの理由から後者を勧めており、正規近似の方法を Cochrane レビューの解析には勧めていません。R ではそれぞれ mada パッケージの reitsma() と lme4 パッケージの glmer() で当てはめられ、例題では結果はほとんど同じになります（発展の\kkref{sec:13-7}）。

\section{\hatten HSROC モデルとの関係}
\label{sec:10-5}

もう 1 つの代表的な方法が、Rutter と Gatsonis（2001）の HSROC モデルです。こちらは、各研究が「陽性と判定する基準の厳しさ（閾値）」と「病気のある人とない人を分ける力（精度）」の 2 つの性質を持ち、その 2 つが研究ごとにばらつくと考えます。閾値の違いは研究を 1 本の曲線の上で動かし、精度の違いは曲線そのものを左上へ寄せたり対角線へ寄せたりします。加えて、閾値が動いたときに精度がどう変わるかを表す形のパラメータがあり、曲線が左右対称かどうかを決めます。

二変量モデルと HSROC モデルは、出発点の考え方は違いますが、共変量を入れない場合には数学的に同じモデルで、片方のパラメータからもう片方のパラメータを計算できます（Harbord 2007）。違いは、結果として何を前面に出すかにあります。二変量モデルは要約点（平均の感度と特異度）とその信頼領域を自然に出し、HSROC モデルは SROC 曲線を自然に出します。

\section{要約点と SROC 曲線のどちらを出すか}
\label{sec:10-6}

SROC 曲線は、ROC 平面（図 10.1）に研究を点で置き、その点の並びに沿って引いた曲線です。カットオフが動いたときに感度と特異度がどう入れ替わるかを、1 本の線で示します。

Cochrane DTA ハンドブックは、要約点と SROC 曲線のどちらを出すかを研究のカットオフで決めるよう整理しています（Macaskill 2023、9.3.1）。組み入れた研究が共通のカットオフを使っているなら、要約点が臨床で意味を持つので、二変量モデルで要約点を推定します。研究が複数のカットオフにまたがり、どのカットオフにも複数の研究があるときは、カットオフごとに同じ解析を繰り返してかまいません。

カットオフが混ざった研究から要約の感度と特異度を出すことは、避けるべきとされています。ハンドブックは、その値が臨床では解釈できず、一般化もできないからだと述べています。例題の PCT には 0.25、0.5、1 ng/mL の研究が混ざっています。これらを平均した感度と特異度は、どのカットオフの値でもなく、PCT を何 ng/mL で陽性としたときの値なのかを答えられません。

カットオフが混ざった全研究には、SROC 曲線を引きます。曲線は、カットオフが動いたときに感度と特異度がどう入れ替わるかの関係を示し、特定のカットオフの値を主張しません。共通のカットオフの研究の要約点と、全研究の SROC 曲線は、両方を示してもかまいません。例題でも両方を出します。要約点はカットオフ 0.5 ng/mL の 6 研究から、SROC 曲線は全 10 研究から出します（\kkref{sec:13-5}、\kkref{sec:13-8}）。カットオフの違いが感度と特異度をどれだけ動かすかは、カットオフを共変量にして調べます（\kkref{sec:14-2}）。

\section{異質性をどう見るか}
\label{sec:10-7}

第 1 部では、異質性を Q、I²、τ² で表しました\kksee{sec:06-1}。検査のメタ解析では、このうち I² を使うことは勧められていません。Cochrane DTA ハンドブックは、感度と特異度それぞれについての等しさの検定と I² の推定を、推奨しないと明記しています（Macaskill 2023、9.3.2）。感度と特異度を別々に計算した I² は、2 つの相関と閾値効果を無視しているからです。

\begin{hattenbox}
そのほかの理由もあります。診断精度研究の感度や特異度は、研究ごとの差が標本誤差よりはるかに大きいことが多く、I² はしばしば 90\% を超えるので、研究どうしを見比べる情報になりません。例えば Wacker らの PCT のメタ解析（2013）では I² は 96\% でした。さらに、二変量モデルに対応する I² の定義はいくつか提案されていますが、まだ定まっていません。例題のカットオフ 0.5 ng/mL の 6 研究を mada で解析すると、提案された方法によって 0\% から 63.7\% まで違う値が出ます。madad() が出力する等しさの検定（\kkref{sec:13-1}）も同じ理由で、出力には出ますが判断には使いません。
\end{hattenbox}

代わりに使うのは、ROC 平面に研究を置いた図と、二変量モデルが推定する研究間の標準偏差、そして 95\% 予測領域です。予測領域は、次の研究がどこに来そうかを示す楕円で、第 1 部の予測区間（\kkref{sec:06-3}）を 2 次元に広げたものです。予測領域が広ければ、要約点がどれだけ確かに推定されていても、実際の現場での精度は大きくばらつくと読みます。ただし予測領域は、研究間の分布が 2 次元の正規分布に従うという仮定に頼っています。研究数が少ないと、この仮定のもとでの推定が不安定で、領域は実際より広くも狭くも出ます（Macaskill 2023、9.3.2、Deeks 2019）。形と大きさは目安として読みます。図を見て、閾値効果で説明できる散らばり（曲線に沿った散らばり）と、説明できない散らばり（曲線から離れる散らばり）を区別し、後者の原因を共変量で探します（\kkref{sec:14-1}）。

\section{小規模研究効果の調べ方}
\label{sec:10-8}

検査のメタ解析では Egger 検定を使わず、調べるなら Deeks の検定を使います（\kkref{sec:15-2}の GRADE でもこの検定を使います）。Egger 検定は\kkref{sec:07-3}で説明しました。理由と Deeks の検定の中身は、発展の\kkref{sec:10-8-1}に書きました。研究が少ないと検出力は低く、有意でなくても偏りがないとは言えません。

\subsection{\hatten Egger 検定を使えない理由と Deeks の検定}
\label{sec:10-8-1}

第 1 部ではファネルプロットと Egger 検定で小規模研究効果を調べました\kksee{sec:07-2}。検査のメタ解析では、Egger 検定を診断オッズ比にそのまま当てはめると、出版バイアスがなくても非対称と判定されやすいことが知られています（Deeks 2005）。診断オッズ比の標準誤差は、2×2 表の 4 つの数から計算され、診断オッズ比の大きさそのものとも結びついているためです。

Deeks らはこの問題を避ける検定を提案しました。縦軸に診断オッズ比の対数、横軸に「有効な標本の大きさ」の平方根の逆数をとって研究を置き、両者の関係を回帰で調べます。有効な標本の大きさは、病気のある人の数とない人の数から計算する量で、標準誤差の代わりに研究の大きさを表します。Cochrane DTA ハンドブックはこの方法を紹介しています。ただし研究数が少ないと検出力が低く、また検査の研究がどのような理由で未発表になるかは介入の試験ほど分かっていないので、検定の結果は慎重に読みます。例題の 10 研究では、この検定で偏りがないと判断できるだけの情報はありません。

一方で、小さい研究ほど精度が高く出るという現象は、検査のメタ解析でも報告されています。Tang らの PCT のメタ解析（2007）では、Sackett と Haynes の分類で phase 3 とされた 4 研究で、標本サイズが小さい研究ほど精度が高く、大きい研究では帰無に近づく関係が見られました（メタ回帰の係数 −0.592、p = 0.017）。検定の有意性よりも、こうした傾向が見えたときに、小さい研究の設計に偏りがないかを QUADAS で確かめることのほうが役に立ちます。

\section{この章のまとめ}
\label{sec:10-9}

検査のメタ解析では、感度と特異度を別々に統合せず、logit の目盛りに直したうえで、両者の相関と研究間のばらつきを同時に扱う二変量モデルか HSROC モデルを使います。共通のカットオフの研究なら二変量モデルで要約点を出し、カットオフが混ざった研究から要約点を出すことは避け、全研究には SROC 曲線を示します。両方を示してもかまいません。異質性は等しさの検定や I² ではなく、ROC 平面の図と予測領域で見ます。予測領域は正規分布の仮定に頼るので、研究が少ないときは目安にとどめます。小規模研究効果は Egger 検定ではなく Deeks の検定で調べ、結果は慎重に読みます。次の第 11 章では、こうした解析に入る前の工程として、問いの立て方と文献の選別を、架空の PCT の抄録で練習します。
